\documentclass[14pt, timesnewroman]{article}
\usepackage[utf8]{inputenc}
\usepackage[russian]{babel}
\usepackage{graphicx} % Required for inserting images

\title{Разработка веб-приложения администрирования групповых занятий}
\author{Сыров Максим Евгеньевич}
\date{Февраль 2025}


\begin{document}
\maketitle 
\section{Предметная область и проблема}

В настоящее время в сфере образования и услуг сложилась необходимость в простом, доступном и развитом инструменте для выполнения организационных задач. Сейчас для организации образовательного процесса преподавателям необходимо произвести много ручных и монотонных действия для контроля посещаемости своих занятий.
\section{Целевая аудитория}
Особенно остро эта проблема возникает у учителей, преподавателей и репетиторов, которым приходится тратить много усилий и времени для сбора данных о своих учениках, выставления им посещаемости, учета заработков с определенных занятий.

Также, с похожей проблемой сталкиваются представители сферы услуг по причине отсутствия единой платформы, где они могут быстро и удобно составить своё расписание, а также предоставить доступ и согласовать мероприятия с своими клиентами, для оптимизации своих временных расходов

\section{Аналоги и конкуренты}
Мое приложение должно выполнять общий необходимый функционал необходимый для решения проблем, описанных выше, а также  не должно иметь их недостатков, препядствующие к достижению итоговой цели применения.

К сожалению, сегодня, большинство преподавателей и учетелей предпочитают использовать редакторы .xslx файлов, такие как Excel и GoogleDocx. В данном случае я объединяю их в один аналог, поскольку эти редакторы совмещают в себе один общий функционал.
Почему его выбирают:
\begin{itemize}
    \item Привычность. Многие используют Excel поскольку привыкли к ячеечному строю и формированию списков своих групп в этом формате.
    \item Доступность. Excel предустановлен на всех Windows устройствах и доступен сразу "из коробки". Пользователю сразу доступен весь функционал.
    \item Удобство хранения. Некоторые xslx редакторы поддерживают облачное хранение, обеспечивая своих пользователем общим доступом с разных устройств, лишая их необходимости носить переносные устройства.
\end{itemize}
Недостатки:
\begin{itemize}
    \item Сбор информации. Пользователю необхожимо вручную вводить информацию по каждому своему ученику на каждое занятие
    \item Дописать еще...
\end{itemize}

Ярепетитор - прямой конкурент моей платформе, который предоставляет общий функционал создания заняий, просмотра расписания, учета заработка, учета оплаченных часов.
Однако эта платформа обладает главным недостатком - односторонност использования. Система предоставляет функционал только для преподавателя. 
Это накладывает следующие сложности:
\begin{enumerate}
    \item Пользователю приходится вручную заносить своих учеников в систему, всю информацию о них
    \item Пользователю приходится самому уведомлять своих учениках о переносе занятия
    \item Ученики пользователя не имеют возможности получить общую информацию о курсах и занятиях преподавателя без предварительной коммуникации. Это пораждает лишнее потраченное время на согласование времени, цены и условия проведения занятий. 
\end{enumerate}


Таким образом, моя система должна обладать следующим функционалом:

Система будет делить функционал на 3 роли: Ученик, Учитель и администратор. (В данном случае администратор включает в себя общий функционал как участника так и ученика для поддержки системы)

Функционал ученика:
\begin{itemize}
    \item Регистрация/Авторизация как ученик
    \item Редактирование своих данных
    \item Просмотр, вступление и покидание групп учителей
    \item Отправка и просмотр своих отправленных запросов на вступление в закрытые группы
    \item Просмотр своего расписания будущих занятий на день/неделю/месяц
    \item Просмотр отдельного занятия
    \item Просмотр своей статистики по прошедшим зантиям по различным параметрам
    \item Возможность принять участие в открытых занятиях учителях (занятиях, не связанные с группами, полезно для тех кто не ведет большое количество учеников)
\end{itemize}

Функционал учителя:
\begin{itemize}
    \item  Регистрация/Авторизация как учитель
    \item  Редактирование своих данных
    \item  Создание, поиск, редактирование, удаление групп с участниками по тематике. В данном случае функционал создания групп учеников по предметному или возрастному признаку. 
    \item  Просмотр запросов от учеников на вступления в группы. В данном случае просмотр идет сгруппированным по группам.
    \item  Принятие и отклоение запросов на вступление в группу
    \item  Выход из группы
    \item  Создание ссылок приглашений в группу
    \item  Создание регулярных занятий внутри одной группы с расписанием и условиями проведения
    \item  Просмотр своего будущего расписания на день/неделю/месяц
    \item  Уведомление участников группы о переносе времени или даты занятия
    \item  Создание, редактирование, поиск, удаление слотов
    \item  Просмотр статистики посещений целой группы или отдельного занятия
    \item  Просмотр своего расписания планируемых занятий
    \item  Ведение посещаемости, оплаты и прочих характеристик на каждом занятии по каждому ученику
\end{itemize}

\section{Задачи}


Задачи:
\begin{enumerate}
    \item Проанализировать предметную область
    \item Исследовать ЦА
    \item Проанализировать конкурентов и аналолги
    \item Определиться с функциональными и нефункциональными требованиями
    \item Спроектировать и описать техническое решение
    \item Спроектировать и разработать систему (веб приложение)
    \item Протестировать разработанную систему на реальных потребителях
    \item Сделать выводы о результатах применения системы и выполненных целях
\end{enumerate}



\section{План по параграфам}

\subsection{I глава}
\begin{itemize}
    \item Анализ предметоной области
    \item Исследование ЦА
    \item Определение целей и задач которые ЦА ставит перед с собой
    \item Исследование существующих средств, выполняющие функции ЦА
    \item Выявление проблем с которыми сталкивается ЦА
\end{itemize}
\subsection{II Глава}
\subsubsection{Проектирование}
\begin{itemize}
    \item Определение функцональных требований
    \item Определение нефункциональных требований
    \item Составление сценариев использования системы
    \item Описание технических требований к системе
\end{itemize}
\subsubsection{Разработка технической части}
\begin{itemize}
    \item Поиск оптимального стека технологий необходимый для выполнения требований
    \item Проектирование базы данных
    \item Описание и разработка необходимых компонентов системы, работающие с каждой таблицей
    \item Разработка необходимых вспомогательных компонентов для работы по HTTP
    \item Юнит-тестирование
\end{itemize}
\subsubsection{Разработка интерфейса}
\begin{itemize}
    \item Из анализа ЦА определение единого стиля и внешнего вида интерфейса ( UI kit)
    \item Разработка внешнего вида всех компонентов, форм и таблиц
    \item Прототипирование
\end{itemize}
\subsubsection{Разработка веб-части}
\begin{itemize}
    \item Поиск оптимального стека технологий необходимый для выполнения требований
    \item Написание API
    \item Разработка UI, обхих компонентов и базовой функциональности
    \item Создание и наполнение страниц
    \item Срщивание с бекендом
    \item Интеграционное тестирование функционала
\end{itemize}

\subsection{III Глава}
\begin{itemize}
    \item Ввод системы в эксплуатацию
    \item Сбор обратной связи и QA
    \item Определения списка результатов по которым можно определить удовлетворенность пользователей
\end{itemize}




\end{document}
